\begin{table}[h!]
\centering
\caption{Key Differences Between Traditional and Lean Manufacturing}
\begin{tabular}{|p{3cm}|p{6cm}|p{6cm}|}
\hline
\textbf{Feature} & \textbf{Lean Manufacturing} & \textbf{Traditional Manufacturing} \\ \hline

\textbf{Production Model} & 
\textbf{Pull System} – Production based on real-time demand & 
\textbf{Push System} – Production based on forecasts \\ \hline

\textbf{Waste Management} & 
Eliminates waste through \textbf{Continuous Improvement} & 
High waste due to excess inventory and overproduction \\ \hline

\textbf{Inventory Management} & 
\textbf{Just-in-Time (JIT)} – Minimum stock levels maintained & 
\textbf{Large inventories} held to prevent shortages \\ \hline

\textbf{Flexibility} & 
High – Easily adapts to demand changes & 
Low – Fixed processes make rapid adjustments difficult \\ \hline

\textbf{Employee Involvement} & 
High – Encourages problem-solving and efficiency improvements & 
Low – Employees follow standardized tasks with limited decision-making \\ \hline

\textbf{Process Optimization} & 
Uses \textbf{Value Stream Mapping} and \textbf{Standardized Work} to refine processes & 
Rigid processes with minimal optimization \\ \hline

\end{tabular}
\end{table}
